\section{模型假设与统一符号}

\subsection{主要假设}

\begin{enumerate}
\item 每个 10 min 区间内，给定功率值可代表该区间平均功率；线路损耗与储能自放电忽略。
\item 微网可以弃光，但题目未要求向外网售电，因此本文只允许 $g_t\ge 0$，不设置上网收益。
\item ``充放电效率为 90\%''解释为充电、放电单向效率均为 $\eta=0.9$；这一口径列为强假设。
\item 问题一明确要求 0:00 与 24:00 电量相同，因此施加终端等式；问题二至四只保持跨日 SOC 连续与安全边界，不额外强制每日回到日初 SOC。
\item 问题三的下调结算主口径为：取消的计划电量不再按原价全额支付，但按该部分 50\% 的时点电价支付违约结算。另一种``计划全额支付后再加 50\% 罚金''的解释单独做敏感性分析。
\item 问题四正式策略不能看到未来真实电价；当前无专门价格预测附件，因此用历史价格构造因果预测，真实当日电价只用于结算。
\end{enumerate}

\subsection{统一变量}

以 $\Delta t=1/6$ h 为时间步长，$t=1,\ldots,144$。定义：$L_t$、$P_t$ 分别为负载与光伏功率（kW）；$p_t$ 为外网电价（元/kWh）；$g_t$ 为外网购电量（kWh）；$c_t,d_t$ 为储能交流侧充、放电量（kWh）；$E_t$ 为区间末储能电量（kWh）。储能安全边界与单步功率约束为
\[
1200\le E_t\le10800,\qquad 0\le c_t,d_t\le 5000\Delta t=833.333.
\]

能量平衡写为
\[
g_t+d_t-c_t\ge (L_t-P_t)\Delta t,
\tag{1}
\]
其中右端为负时允许弃光。储能状态转移为
\[
E_t=E_{t-1}+\eta c_t-\frac{d_t}{\eta},\qquad \eta=0.9.
\tag{2}
\]

由于外网电价为正且往返储能存在损耗，在当前 LP 最优解中没有同一时段同时充放电；确认实验中同时充放电次数为 0，因此无需为了``互斥''额外引入二进制变量。
